\section{Discusión}

\subsection{7.1. Ventajas y desventajas comparativas}

Resulta revelador observar cómo ninguna arquitectura emerge como claramente superior en todos los escenarios. Los combustores anulares demuestran ventajas sustanciales en eficiencia de combustión, uniformidad térmica y reducción de emisiones, alineándose bien con los objetivos ambientales más exigentes. 

Liu et al. refuerzan esta visión al mostrar trade-offs similares cuando se incorporan combustibles alternativos, donde la mejor mezcla del anular se traduce en menor formación de soot y mejor tolerancia a variaciones de combustible \cite{Liu2024}. 

Sin embargo, los can-anulares conservan fortalezas operativas difíciles de ignorar: mayor robustez en transitorios, facilidad de mantenimiento y una madurez tecnológica que reduce riesgos en certificación. Particularmente llamativo es el hecho de que estas ventajas prácticas sigan pesando fuertemente en decisiones de flota actual, a pesar de las brechas en desempeño ambiental.

Aunque los enfoques anulares tienden a ofrecer mejores números en papel, los can-anulares siguen representando una solución más pragmática para fabricantes que priorizan tiempo al mercado y costos de ciclo de vida.

\subsection{7.2. Implicaciones para diseño de motores futuros}

Uno de los desafíos persistentes que surge al proyectar estas tecnologías hacia la próxima generación de turbofans radica en cómo integrarlas en ciclos de muy alto bypass y presiones de operación crecientes. Los resultados sugieren que los combustores anulares se posicionan como la opción natural para motores destinados a cumplir metas net-zero, especialmente cuando se combinan con SAF o hidrógeno.

Moreno et al. aportan insights valiosos desde sus simulaciones CFD, demostrando que las geometrías continuas facilitan la implementación de inyección multipunto y control activo de dilución, elementos clave para combustión lean premixed estable \cite{Moreno2024}. 

Lejos de ser un asunto resuelto, esta transición exigirá avances paralelos en materiales cerámicos de matriz (CMC) para paredes expuestas a mayores temperaturas y en sistemas de fabricación aditiva que reduzcan costos de los anulares. Los diseñadores deberán además considerar la integración con sistemas de control del motor completo, donde la uniformidad del anular podría simplificar algoritmos de salud y pronóstico.

\subsection{7.3. Limitaciones del estudio}

No obstante, cabe matizar que el presente análisis comparativo, aunque riguroso en su marco metodológico, presenta limitaciones inherentes que merecen reconocimiento explícito. Las simulaciones CFD, pese a su alta resolución, no sustituyen pruebas en banco a presión completa ni ensayos en vuelo. 

Liu complementa esta precaución al resaltar en estudios similares cómo las interacciones entre combustor y turbina pueden modificar sensiblemente los perfiles reales de temperatura \cite{Liu2024}. 

Además, el análisis se centró principalmente en condiciones estacionarias representativas, dejando espacio para mayor profundidad en comportamientos transitorios y efectos de envejecimiento. La disponibilidad limitada de datos experimentales públicos para configuraciones exactas de motores comerciales actuales obligó a ciertas simplificaciones en los modelos de combustión y propiedades de materiales.

Estas limitaciones, lejos de invalidar las tendencias observadas, subrayan la necesidad de validación experimental colaborativa entre industria y academia. Futuros trabajos deberían incorporar ensayos en instalaciones de alta presión y explorar configuraciones híbridas o activamente controladas que combinen lo mejor de ambas arquitecturas.